\documentclass[11pt]{article}

\usepackage[a4paper,margin=1in]{geometry}
\usepackage{amsmath,amssymb,amsthm,mathtools}
\usepackage{hyperref}
\hypersetup{colorlinks=true,linkcolor=blue,citecolor=blue,urlcolor=blue}

\newtheorem{definition}{Definition}
\newtheorem{theorem}{Theorem}
\newtheorem{proposition}{Proposition}
\newtheorem{remark}{Remark}
\newtheorem{example}{Example}

\newcommand{\NN}{\mathbb{N}}
\newcommand{\RR}{\mathbb{R}}
\newcommand{\CC}{\mathbb{C}}

\title{Compact and Finite-Rank Operators (Quick Notes)}
\author{}
\date{}

\begin{document}
\maketitle

\section*{Setting}
Let $X,Y$ be normed (often Banach) spaces over $\RR$ or $\CC$, and let $\mathcal{B}(X,Y)$ denote the space of bounded linear operators $T:X\to Y$ with the operator norm
\[
\|T\|=\sup_{\|x\|\le 1}\|Tx\|.
\]

\section{Finite-Rank Operators}
\begin{definition}[Finite rank]
A bounded linear operator $T:X\to Y$ has \emph{finite rank} if its range $T(X)$ is a finite-dimensional subspace of $Y$. The \emph{rank} of $T$ is $\dim T(X)$.
\end{definition}

\begin{example}
\leavevmode
\begin{itemize}
  \item \textbf{Rank one.} If $f\in X'$ and $y_0\in Y$, the map $T(x)=f(x)\,y_0$ has range $\mathrm{span}\{y_0\}$, hence rank $1$.
  \item \textbf{Finite sum of rank ones.} If $T(x)=\sum_{k=1}^n f_k(x)\,y_k$ with $f_k\in X'$ and $y_k\in Y$, then $T(X)\subseteq \mathrm{span}\{y_1,\dots,y_n\}$, so $T$ has finite rank.
  \item \textbf{Coordinate truncation on $\ell^2$.} $T(x_1,x_2,\dots)=(x_1,\dots,x_n,0,0,\dots)$ has rank $n$.
\end{itemize}
\end{example}

\section{Compact Operators}
\begin{definition}[Compact operator]
A bounded linear operator $T:X\to Y$ is \emph{compact} if it maps bounded sets into relatively compact sets. Equivalently: for every bounded sequence $(x_m)\subset X$ there is a subsequence $(x_{m_j})$ such that $(Tx_{m_j})$ converges in $Y$.
\end{definition}

\begin{example}
\leavevmode
\begin{itemize}
  \item \textbf{Finite rank $\Rightarrow$ compact.} In finite-dimensional spaces, bounded sets are relatively compact (Heine--Borel). Thus every finite-rank operator is compact.
  \item \textbf{Diagonal operator on $\ell^2$.} Define $T:\ell^2\to\ell^2$ by
  \[
    T(x_1,x_2,\dots)=\bigl(\tfrac{x_1}{1},\tfrac{x_2}{2},\tfrac{x_3}{3},\dots\bigr).
  \]
  Then $T$ is compact: the singular values $1/n\to 0$, so $T$ is the norm-limit of finite-rank truncations.
  \item \textbf{Integral operator on $C([0,1])$.} If $K\in C([0,1]^2)$ and
  \[
    (Tf)(x)=\int_0^1 K(x,y)f(y)\,dy,
  \]
  then $T:C([0,1])\to C([0,1])$ is compact (Arzel\`a--Ascoli).
\end{itemize}
\end{example}

\begin{remark}[Non-examples]
\leavevmode
\begin{itemize}
  \item The identity $I_X$ on an infinite-dimensional Banach space $X$ is \emph{not} compact: the unit ball is bounded but not relatively compact.
  \item The inclusion $\iota:\ell^2\hookrightarrow \ell^\infty$ is bounded but \emph{not} compact: $(e_n)$ lies in the $\ell^2$ unit ball and $(\iota e_n)$ has no convergent subsequence in $\ell^\infty$.
\end{itemize}
\end{remark}

\section{Basic Theorems}
\begin{proposition}\label{prop:fr-compact}
Every finite-rank operator is compact.
\end{proposition}
\begin{proof}[Proof sketch]
If $\dim T(X)<\infty$, then $T(B)$ is bounded in a finite-dimensional space, hence relatively compact. Apply the definition.
\end{proof}

\begin{theorem}[Closedness and ideal properties]
Let $X,Y,Z$ be Banach spaces. The set $\mathcal{K}(X,Y)$ of compact operators is a closed linear subspace of $\mathcal{B}(X,Y)$. Moreover, if $A\in\mathcal{B}(Y,Z)$ and $B\in\mathcal{B}(X,Y)$, then 
\[
A\circ T\circ B\in \mathcal{K}(X,Z)\quad\text{whenever }T\in\mathcal{K}(Y).
\]
Thus $\mathcal{K}$ is a two-sided operator ideal.
\end{theorem}

\begin{proposition}[Riesz lemma consequence]
If $X$ is infinite-dimensional and $T\in\mathcal{B}(X)$ is compact and injective, then $T$ cannot have a bounded inverse on $T(X)$; in particular, the identity operator is not compact.
\end{proposition}

\begin{theorem}[Sequential characterization]
$T:X\to Y$ is compact $\iff$ for every bounded sequence $(x_m)\subset X$ there exists a subsequence $(x_{m_j})$ such that $(Tx_{m_j})$ converges in $Y$.
\end{theorem}

\begin{proposition}[Weak-to-strong property]
If $X$ is a Banach space and $T\in\mathcal{K}(X,Y)$, then whenever $x_n\rightharpoonup 0$ weakly in $X$ and $(x_n)$ is bounded, we have $\|Tx_n\|\to 0$.
\end{proposition}

\section{Relations and Remarks}
\begin{itemize}
  \item Finite-rank $\subset$ compact $\subset$ bounded, with both inclusions proper in infinite dimensions.
  \item Compact operators generalize ``finite-matrix'' behavior: spectral theory for compact operators on infinite-dimensional Hilbert spaces mimics the finite-dimensional case (e.g.\ nonzero spectrum is discrete with $0$ as the only possible accumulation point).
  \item In many classical spaces (e.g.\ Hilbert spaces, $C(K)$ with $K$ compact metric), compact operators can be approximated in norm by finite-rank operators; in full generality this relies on approximation properties of the spaces involved.
\end{itemize}

\section*{Mini Checklist / Examples at a Glance}
\begin{itemize}
  \item Rank-one: $T(x)=f(x)y_0$ (finite-rank $\Rightarrow$ compact).
  \item Diagonal on $\ell^2$ with entries $s_n\to 0$ is compact; if $s_n\nrightarrow 0$, it is not compact.
  \item Integral operator with continuous kernel on $C([0,1])$: compact (Arzel\`a--Ascoli).
  \item Identity on infinite-dimensional space: not compact.
\end{itemize}

\end{document}
